\scanpage{110}
\mspage{5}
% archive foliation: 109

La donnée d'un point $x$ \inserted{$\ne0$,} de $E(k)$ équivaut à celle de son image
\unsure{dans} $\mathbb P^1_k$, soit \deleted{$\unsure{\wp(x)}$}
\[
(38)\qquad \wp(x)=x_0\in k,
\]
plus la donnée d'une des deux solutions \inserted{$x'$} de l'équation
(\uline{deux} tout au moins si le 2\textsuperscript{e} mb est $\ne0$ i.e.
si $x\notin J={}_2E^*$ i.e. $x_0\notin D_J\subset\Sigma_J$)
\[
(39)\qquad X'^2=x_0^3+a_1x_0^2+a_2x_0+a_3,
\]
donnée en termes de $x$ par
\[
(40)\qquad x'=\wp'(x).
\]

\uline{Proposition 3} Soient $x,y,z\in E(k)\setminus\{0\}$,
\inserted{\deleted{pas} tous \deleted{égaux} \uline{distincts},} alors on a
\[
(41)\qquad x+y+z=0
\]
SSS on a \deleted{\unsure{ou bien} $\wp(x)=\wp(y)=\wp(z)$ \unsure{ou}}
\[
(42)\qquad \det\begin{pmatrix}
1&\wp(x)&\wp'(x)\\
1&\wp(y)&\wp'(y)\\
1&\wp(z)&\wp'(z)
\end{pmatrix}=0.
\]
% margin position: left beside (41)--(42)
\margin{\uline{NB} Si deux parmi les $x,y,z$ sont égaux,
(42) est tjrs vérifié, sans qu'on ait \unsure{nécessairement} $x+y+z=0$ !}
% margin position: right beside (41)--(42)
\margin{Cela implique qu'on a
$\wp(x)\ne\wp(y)\ne\wp(z)\ne\wp(x)$
(car si on avait $\wp(x)=\wp(y)$ on aurait $x=y$ ou $x=-y$, d'où $z=0$)
\unsure{Donc le} \unsure{cas est|l'autre} \unsure{exclu}.}
% margin position: left below (42)
\margin{\deleted{La condition $\unsure{\to|\Rightarrow}$ ici,
(et non $\unsure{\leftarrow|\Leftarrow}$) signifie
$\exists a\in E(k)$, avec deux \deleted{\unsure{des}} \unsure{des}
\inserted{trois} points $x,y,z$ opposés ; le troisième $=a$,
d'où $x+y+z=a$, donc $a=0$, ce qui est absurde puisque $x,y,z\ne0$.}}

En effet, la définition de la loi de multiplic. sur $E$ nous donne que si
$x,y,z\in E(k)$ sont \unsure{arbitraires}, alors \deleted{$x+y+z=0$}
\[
x+y+z=0\quad\Longleftrightarrow\quad
(\{x\}-\{0\})+(\{y\}-\{0\})+(\{z\}-\{0\})\overset{\mathrm{lin.}}\sim0
\]
\[
\text{i.e.}\qquad \{x\}+\{y\}+\{z\}\overset{\mathrm{lin.}}\sim3\{0\}
\]
Si $x,y,z$ sont tous $\ne0$, cela signifie qu'il existe une fonction rat. sur $E$,
ayant $0$ comme pôle d'ordre 3, pas d'autre pôle, (donc $f\in\mathcal F_3$),
et $\{x\}+\{y\}+\{z\}$ comme \deleted{pôle} diviseur des zéros.
